% titel: Blatt 4, A6: Elektronenkonfiguration und Quantenzahlen
% skript: Mortimer Kap. 6.7 (Quantenzahlen, Elektronenkonfiguration)

% Teil I: Werkzeuge, kein Lösungsweg.

Elektronen eines Atoms werden durch Quantenzahlen beschrieben; die
Elektronenkonfiguration sortiert die Elektronen nach Orbitalen (s, p, d, f)
und Schalen (n). Der Zusammenhang zur Periodentafel ist direkt: die
Hauptquantenzahl $n$ der Valenzschale entspricht der Periodennummer.

\begin{werkzeugkasten}[Quantenzahlen]
\begin{itemize}
  \item \textbf{Hauptquantenzahl} $n = 1,2,3\,\dots$: Schale/Energie;
        Periodennummer des Elements im PSE.
  \item \textbf{Nebenquantenzahl} $l = 0,\dots,n-1$: Orbitalform;
        $l=0$ ($s$), $1$ ($p$), $2$ ($d$), $3$ ($f$).
  \item \textbf{Magnetische Quantenzahl} $m_l = -l,\dots,+l$:
        Orientierung; $2l+1$ Werte pro $l$.
  \item \textbf{Spin} $m_s = \pm \tfrac12$.
\end{itemize}
Konfiguration bestimmt man nach Aufbauprinzip (Füllung $1s < 2s < 2p <
3s < 3p < 4s < 3d < 4p \,\dots$) und zählt die Elektronen bis $Z$.
Stabile Edelgas-Rümpfe ($\mathrm{[He]}$, $\mathrm{[Ar]}$, $\mathrm{[Kr]}$,
$\mathrm{[Xe]}$) kann man abkürzen.

\textit{Skript:} Mortimer Kap. 6.7 (Quantenzahlen, Aufbauprinzip), Übung 4/6.
\end{werkzeugkasten}

\begin{achtung}[Darauf achten]
\begin{itemize}
  \item Zu jedem $l$ gehören genau $2l+1$ Werte von $m_l$; die Anzahl der
        Orbitale in s, p, d, f ist 1, 3, 5, 7.
  \item 3d wird erst \emph{nach} 4s gefüllt (also für $Z\ge 21$); nicht
        $3d$ vor $4s$ schreiben.
  \item Die \emph{äusserste} Schale bestimmt die Periodennummer (z.\,B.
        Pb: Perioden 6 $\Rightarrow$ Valenzschale $n=6$).
  \item Für $d$- und $f$-Block-Elemente Aufbau- und
        Halbbesetzung-Regeln beachten; hier genügt das Grundprinzip.
\end{itemize}
\end{achtung}
